\documentclass[11pt,a4paper]{article}

% --- Pacchetti di base ---
\usepackage[utf8]{inputenc}
\usepackage[T1]{fontenc}
\usepackage[italian]{babel}
\usepackage[margin=2.5cm]{geometry}
\usepackage{lmodern} % Font leggibile e standard

% --- Pacchetti per tabelle e grafica ---
\usepackage{booktabs}
\usepackage{tabularx}
\usepackage{graphicx}
\usepackage{xcolor}

% --- Pacchetti per layout e link ---
\usepackage{hyperref}
\usepackage{enumitem}

% Definizione colore istituzionale (Blu Bicocca)
\definecolor{bicoccablue}{RGB}{0, 84, 150}

\hypersetup{
	colorlinks=true,
	linkcolor=bicoccablue,
	urlcolor=bicoccablue,
	pdftitle={Rapporto di Follow-up SDL - Bicocca Robotic Chemical Lab},
}

% --- Informazioni Titolo ---
\title{\vspace{-1.5cm}\textbf{\huge Rapporto di Follow-up}\\ \Large Stato di Avanzamento e Risultati del Laboratorio Chimico Robotico Autonomo (SDL)}
\author{Bicocca Robotic Chemical Lab}
\date{Aprile 2026}

\begin{document}
	
	\maketitle
	
	% --- EXECUTIVE SUMMARY ---

	\section*{Executive Summary}
	Questa fase del progetto segna il passaggio critico dal design concettuale all’integrazione funzionale e alla validazione sperimentale dell'infrastruttura robotica. In un panorama scientifico dove l’Intelligenza Artificiale accelera esponenzialmente la scoperta di nuovi materiali, emerge chiaramente il cosiddetto "AI-Experimental Gap": la capacità di generare candidati teorici supera di ordini di grandezza la velocità della validazione manuale. L'obiettivo centrale è l'implementazione della Flexible Automation, un paradigma che supera la rigidità dei sistemi industriali tradizionali per rispondere alla natura dinamica della ricerca chimica, dove i protocolli evolvono quotidianamente.
	I tre contributi strategici di questa fase includono:
	\begin{itemize}[]
	\item \textbf{Architettura Modulare:} Una scomposizione del laboratorio in workstation indipendenti e funzionali (Liquid Handling, Analytical, Manual), orchestrate da un sistema di logistica autonoma.
	\item \textbf{Integrazione SiLA2:} L'adozione di uno standard di comunicazione aperto per garantire l’interoperabilità ed eliminare il rischio di "vendor lock-in", assicurando la scalabilità del sistema.
	\item \textbf{Validazione End-to-End:} La verifica sperimentale su un processo reale di sintesi di idrogeli, dimostrando la capacità dell’infrastruttura di gestire l’intero ciclo di vita del campione con tracciabilità totale.
	\end{itemize}
	
	\noindent L'approccio proposto trasforma il laboratorio in un ecosistema integrato che si pone come obiettivo la transizione verso l'automazione flessibile .
	
	
	\newpage
	\tableofcontents
	\newpage
	
	\section{Architettura di Sistema e Standardizzazione SiLA2}
	L’adozione di standard aperti come SiLA2 (Standardization in Laboratory Automation) è un imperativo strategico per evitare la frammentazione tecnologica. In un ecosistema di ricerca, la capacità di integrare strumenti di produttori diversi senza riscrivere la logica core del sistema è il principale driver di efficienza a lungo termine.
	
	Il sistema si basa su una Dual Interface Architecture. Ogni server espone sia le "Native SiLA2 Features" (per operazioni specifiche dello strumento) sia il protocollo SiLA2Common. Quest'ultimo funge da strato di astrazione universale che abilita il "Plug-and-Play": l'orchestratore scopre le capacità dei dispositivi a runtime, permettendo scalabilità illimitata. L'aggiunta di nuovi strumenti non richiede riconfigurazioni del software di base; ogni nuova risorsa viene integrata nel workflow non appena connessa alla rete.
	
	 Inoltre, per garantire la longevità dei dati e il loro utilizzo in futuri training di IA, il sistema supporta l'esportazione dei risultati in formato AnIML (Analytical Information Markup Language).

	
	\begin{table}[h!]
		\centering
		\renewcommand{\arraystretch}{1.3}
		\begin{tabularx}{\textwidth}{@{} l X @{}}
			\toprule
			\textbf{Componente Hardware} & \textbf{Funzione Principale} \\ \midrule
			Opentrons Flex & Liquid Handling (Pipetting, Diluizioni) \\
			Tecan Infinite 200 PRO &  Analytical Measurement (Assorbanza, Fluorescenza, Luminescenza) \\
			GoFaGo (Mobile Robot) & Logistics (Trasporto autonomo campioni) \\
			Manual Station & Human-in-the-loop (Task complessi) \\ \bottomrule
		\end{tabularx}
		\caption{Mappatura dell'integrazione hardware tramite standard SiLA2}
	\end{table}
	

	\section{Logistica Autonoma}
	Il Mobile Manipulator (MoMa) GoFaGo funge da tessuto connettivo del laboratorio, eliminando la necessità di nastri trasportatori fissi e permettendo una riconfigurazione totale del layout. Il sistema integra la base omnidirezionale Robotnik RB-Kairos+ con il braccio collaborativo ABB GoFa CRB 15000.
	
	Per superare i limiti di precisione della navigazione SLAM, insufficiente per la manipolazione di multiwell plate, abbiamo implementato un sistema di fine-positioning visuale (fine al millimetro). basato sul rilevamento di marker ArUco posti su ogni workstation, il robot compensa l'errore della base ricalibrando la traiettoria del braccio in tempo reale. La facilità d'uso è garantita dalla piattaforma Exsensia, che permette l'apprendimento di nuovi task tramite kinesthetic demonstration (LfD - Learning from Demonstration), dove l'operatore guida fisicamente il braccio senza scrivere codice.
	
	Il trasferimento di una piastra si articola in cinque fasi:
	\begin{enumerate}
		\item \textbf{Navigazione:} Spostamento autonomo verso la stazione di pick tramite SLAM.
		\item \textbf{Localizzazione Marker:} Acquisizione del marker ArUco per definire il sistema di riferimento locale.
		\item \textbf{Pick:} Esecuzione della traiettoria di prelievo (calcolata rispetto al marker).
		\item \textbf{Retraction:} Ritiro dell'arm in configurazione di trasporto sicura entro il footprint della base.
		\item \textbf{Navigazione:} Spostamento autonomo verso la stazione target tramite SLAM.
		\item \textbf{Place:} Posizionamento millimetrico della piastra nella destinazione target.
	\end{enumerate}
 
	
	\section{Automazione dei Workflow}	
	I workflow sono definiti come una sequenza di step, sequenziali o paralleli, che rappresentano la logica dell'esperimento che si vuole eseguire.	
	La democratizzazione della configurazione per il personale non esperto di coding è stata raggiunta tramite un No-code Visual Designer. Questa interfaccia a blocchi permette ai chimici di tradurre la propria expertise di dominio in workflow complessi senza barriere tecniche.
	
	Il valore strategico risiede nell'Hardware Abstraction Layer (HAL) applicato alle ricette di liquid handling. L'HAL separa la logica sperimentale (la "ricetta" con volumi e sequenze) dalla configurazione fisica del deck (dove si trovano gli oggetti). Questo permette la process portability: la stessa ricetta può essere eseguita su diverse configurazioni hardware senza modifiche, riducendo drasticamente i tempi di setup per nuovi esperimenti.
	
	
	\section{Analisi delle Performance}
	La validazione condotta presso il JOiiNT Lab su un processo di Sintesi di Idrogeli ha fornito dati cruciali sulla maturità del sistema. La robustezza software è stata confermata da test di Fault Injection, dimostrando che il sistema gestisce correttamente disconnessioni o errori hardware senza perdita di integrità dei dati.
	
	L'analisi tecnica della manipolazione ha rivelato un'importante asimmetria della tolleranza. Mentre l'asse longitudinale (Y) offre un margine ampio (da -5 a +20 cm), l'asse laterale (X) è limitato a ±4 cm. Questo accade perché l'errore laterale "consuma" il margine cinematico del braccio, spingendolo verso singolarità che ne impediscono il movimento fluido.
	
	 È stato inoltre identificato il "Forward-push effect" della pinza OnRobot RG6: la cinematica di chiusura sposta la piastra di alcuni millimetri in avanti, un fattore che richiede una compensazione software specifica o l'adozione di gripper a cinematica parallela.
	 
	 		
	
	\section{Conclusioni}
	Il sistema ha un'infrastruttura workstation-centric basata su standard SiLA2 è robusta e validata. Dimostrando che l'integrazione di standard aperti e manipolazione mobile è la via maestra per scalare la ricerca chimica automatizzata.
	
	Le aree di miglioramento prioritario per la prossima fase includono:
	\begin{itemize}[]
		\item \textbf{Hardware Refining:} Adozione di marker ad alta precisione e ricalibrazione sistematica delle camere per abbattere l'errore di posizionamento residuo.
		\item \textbf{Manipolazione Avanzata:} Passaggio a pinze a cinematica parallela per eliminare l'effetto push e aumentare il rate di successo nelle workstation più vincolate.
		\item \textbf{AI Integration:}Implementazione di loop decisionali basati sui dati raccolti (AnIML) per trasformare il sistema in un Self-Driving Lab capace di ottimizzare autonomamente le formulazioni chimiche.	
		\end{itemize}

	
\end{document}